\documentclass{article}
\usepackage{spconf,amsmath,graphicx,booktabs}
\usepackage{url}

\title{SPEECHMASTER: CONFIDENCE-ROUTED SELF-SUPERVISED ASR WITH DISCRETE UNIT BUDGETING}
\name{Student Name$^{\star}$}
\address{$^{\star}$Shanghai Jiao Tong University}

\begin{document}
\maketitle

\begin{abstract}
Self-supervised speech models are highly effective for low-resource ASR, yet
using the strongest model on every utterance wastes computation and hides the
representation/cost tradeoff. We present SpeechMaster, a budget-aware ASR
framework that composes heterogeneous SSL representations instead of selecting a
single encoder. SpeechMaster first decodes with a fast Wav2Vec2 CTC branch,
computes an utterance-level CTC uncertainty score from non-blank posterior
peakiness and entropy, and routes only uncertain utterances to a stronger HuBERT
branch. A separate HuBERT discrete-unit auditor measures token rate and bitrate,
linking recognition accuracy to representation compactness. On LibriSpeech clean
validation slices, routing only 25\% of utterances reduces WER from 2.47\% to
1.88\%, while the full HuBERT branch reaches 1.58\%. An oracle route obtains
1.19\% WER, confirming branch complementarity.
\end{abstract}

\begin{keywords}
ASR, self-supervised speech, uncertainty routing, HuBERT, Wav2Vec2
\end{keywords}

\section{Introduction}

Low-resource ASR is limited by the amount of transcribed speech, while modern
self-supervised learning (SSL) encoders reduce this dependence by pretraining on
unlabeled audio. Wav2Vec2, HuBERT, and WavLM learn complementary speech
representations: contrastive contextual features, masked hidden-unit prediction,
and denoising-enhanced masked prediction. A common project design is to choose
one SSL encoder and report WER. This leaves two practical questions unanswered:
when is a lightweight SSL recognizer already sufficient, and when should a
stronger recognizer be invoked?

We propose SpeechMaster, a composed ASR framework for this setting. Its key
idea is to treat SSL recognizers as a routed ensemble under a compute budget.
The fast branch gives a first transcription and a CTC confidence estimate. The
router sends only low-confidence utterances to a stronger SSL branch. In
parallel, a discrete-unit auditor clusters HuBERT hidden states and reports
bitrate, so the same system studies accuracy, speed, and representation
compression. The contributions are: (1) a confidence-routed SSL-ASR framework,
(2) a CTC uncertainty score requiring no extra labels or training, (3) a
discrete-unit budget auditor, and (4) reproducible LibriSpeech experiments with
WER, CER, RTF, token rate, and bitrate.

\section{SpeechMaster Method}

\subsection{Two-Branch SSL Recognizer}

Given waveform $x$, a pretrained SSL encoder $f_\theta$ produces hidden states
$H=f_\theta(x)$. A CTC projection maps frames to vocabulary logits:
\begin{equation}
    p_\theta(y_t \mid x) =
    \mathrm{softmax}(W H_t + b).
\end{equation}
SpeechMaster uses a fast Wav2Vec2 branch $F$ and a stronger HuBERT branch $S$.
The fast branch is always evaluated. The strong branch is evaluated only if the
router decides that the fast branch is uncertain.

\subsection{CTC Uncertainty Router}

For each non-blank frame of the fast branch, SpeechMaster measures posterior
peakiness and entropy. Let $q_t$ be the CTC posterior and $\mathcal{T}$ be the
frames whose argmax is not blank. The confidence score is
\begin{equation}
    c(x)=\frac{1}{|\mathcal{T}|}\sum_{t\in\mathcal{T}}\max_k q_t(k)
    - \lambda \frac{1}{|\mathcal{T}|}\sum_{t\in\mathcal{T}}H(q_t),
\end{equation}
with $\lambda=0.03$. For a routing budget $B$, utterances with the lowest
$c(x)$ are sent to HuBERT. This gives a tunable continuum between fast-only
decoding and strong-only decoding:
\begin{equation}
    \hat{y}(x)=
    \begin{cases}
    S(x), & x \in \mathrm{BottomB}(c),\\
    F(x), & \mathrm{otherwise}.
    \end{cases}
\end{equation}

\subsection{Layer and Unit Auditors}

SpeechMaster includes two auditors for analysis. The layer auditor decodes
Wav2Vec2 final states and a last-four-layer average through the same CTC head,
testing whether the projection is calibrated to a particular depth. The unit
auditor clusters HuBERT hidden states with MiniBatchKMeans and computes token
rate and bitrate:

\begin{equation}
    \mathrm{bitrate}=r_{\mathrm{tok}}\log_2K.
\end{equation}
Consecutive duplicate units are collapsed to estimate a compact content stream.

\section{Experimental Setup}

\subsection{Data and Models}

Experiments use LibriSpeech clean validation slices. The repository also
provides a low-resource fine-tuning script for train-clean-100 subsets. The
fast branch is Wav2Vec2 base fine-tuned on LibriSpeech; the strong branch is
HuBERT large fine-tuned on LibriSpeech; the unit auditor uses HuBERT base.
Audio is resampled to 16 kHz.

\subsection{Metrics}

ASR quality is measured by WER and CER after LibriSpeech-style text
normalization. Inference efficiency is measured by real-time factor (RTF).
Discrete units are evaluated by codebook size, token rate, deduplicated token
rate, bitrate, and deduplicated bitrate. These metrics correspond to the
assignment requirements and make the recognition/compression tradeoff explicit.

\section{Results}

Table~\ref{tab:speechmaster} is the main SpeechMaster result. Routing only the
lowest-confidence 25\% of utterances recovers most of the gain over Wav2Vec2,
reducing WER by 24.1\% relative while adding a small amount of cascade cost.
The oracle row shows the complementarity ceiling if routing decisions were
perfect.

\begin{table}[t]
\centering
\caption{SpeechMaster budgeted routing results.}
\label{tab:speechmaster}
\small
\resizebox{\columnwidth}{!}{%
\IfFileExists{../results/tables/speechmaster_metrics.tex}{\begin{tabular}{llrrl}
\toprule
Route & Routed utt. & WER & CER & RTF \\
\midrule
0\% & 0 & 2.474 & 0.945 & 0.002 \\
25\% & 32 & 1.877 & 0.633 & 0.003 \\
50\% & 64 & 1.792 & 0.603 & 0.003 \\
75\% & 96 & 1.706 & 0.553 & 0.003 \\
100\% & 128 & 1.578 & 0.493 & 0.004 \\
Oracle & - & 1.195 & 0.382 & - \\
\bottomrule
\end{tabular}
}{
\begin{tabular}{lccc}
\toprule
Route \% & WER(\%) & CER(\%) & RTF\\
\midrule
0 & -- & -- & --\\
25 & -- & -- & --\\
100 & -- & -- & --\\
\bottomrule
\end{tabular}}}
\end{table}

Table~\ref{tab:asr} reports ASR performance. Lower WER/CER is better. The
layer-fusion row tests whether averaging SSL layers improves the CTC decision
surface without training a new decoder.

\begin{table}[t]
\centering
\caption{ASR results on LibriSpeech clean validation slices.}
\label{tab:asr}
\small
\resizebox{\columnwidth}{!}{%
\IfFileExists{../results/tables/asr_metrics.tex}{\begin{tabular}{llrrrr}
\toprule
System & Rep. & N & WER & CER & RTF \\
\midrule
wav2vec2-base-960h & final & 128 & 2.474 & 0.945 & 0.003 \\
wav2vec2-base-960h & last4 & 128 & 2.858 & 1.026 & 0.003 \\
hubert-large-ls960-ft & final & 128 & 1.578 & 0.493 & 0.005 \\
\bottomrule
\end{tabular}
}{
\begin{tabular}{lccc}
\toprule
System & WER(\%) & CER(\%) & RTF\\
\midrule
Wav2Vec2 final & -- & -- & --\\
Wav2Vec2 last4 & -- & -- & --\\
HuBERT final & -- & -- & --\\
\bottomrule
\end{tabular}}}
\end{table}

Table~\ref{tab:units} summarizes discrete-unit compression. Increasing $K$
raises the nominal bitrate, while deduplication reduces token rate by merging
stable frame-level assignments. This makes it possible to compare a continuous
CTC system against compact speech-token representations.

\begin{table}[t]
\centering
\caption{HuBERT discrete-unit statistics.}
\label{tab:units}
\small
\resizebox{\columnwidth}{!}{%
\IfFileExists{../results/tables/unit_metrics.tex}{\begin{tabular}{lrrrrr}
\toprule
System & K & Tok/s & Dedup tok/s & Bit/s & Dedup bit/s \\
\midrule
hubert-base-ls960 & 50 & 49.86 & 25.42 & 281.41 & 143.45 \\
hubert-base-ls960 & 100 & 49.86 & 28.15 & 331.28 & 187.05 \\
hubert-base-ls960 & 200 & 49.86 & 29.51 & 381.14 & 225.58 \\
hubert-base-ls960 & 500 & 49.86 & 33.31 & 447.05 & 298.68 \\
\bottomrule
\end{tabular}
}{
\begin{tabular}{lccc}
\toprule
$K$ & Token/s & Dedup token/s & Bit/s\\
\midrule
50 & -- & -- & --\\
100 & -- & -- & --\\
200 & -- & -- & --\\
500 & -- & -- & --\\
\bottomrule
\end{tabular}}}
\end{table}

\subsection{Ablation and Error Analysis}

The ablations cover model family, layer fusion, routing budget, and codebook
size. HuBERT has lower WER than Wav2Vec2, but the oracle route is better than
either individual branch, indicating non-identical errors. The last-four-layer
average is worse than the final Wav2Vec2 layer, suggesting the pretrained CTC
head is calibrated to final-layer content. Increasing codebook size raises
bitrate while reducing clustering distortion, so compact speech units should be
chosen according to the downstream cost budget.

\section{Conclusion}

SpeechMaster combines fast and strong SSL recognizers through CTC uncertainty
routing, then audits the representation cost with discrete HuBERT units. The
results show that a small routed fraction gives a large WER reduction, while
the oracle route reveals remaining complementarity. This supports the central
claim: low-resource SSL-ASR should optimize accuracy, inference budget, and
representation bitrate jointly instead of treating encoder selection as a
single-model decision.

\section{References}

\begin{thebibliography}{9}
\bibitem{librispeech}
V. Panayotov, G. Chen, D. Povey, and S. Khudanpur, ``Librispeech: An ASR
corpus based on public domain audio books,'' in \emph{Proc. ICASSP}, 2015.

\bibitem{wav2vec2}
A. Baevski, H. Zhou, A. Mohamed, and M. Auli, ``wav2vec 2.0: A framework for
self-supervised learning of speech representations,'' in \emph{NeurIPS}, 2020.

\bibitem{hubert}
W.-N. Hsu et al., ``HuBERT: Self-supervised speech representation learning by
masked prediction of hidden units,'' \emph{IEEE/ACM TASLP}, 2021.

\bibitem{wavlm}
S. Chen et al., ``WavLM: Large-scale self-supervised pre-training for full stack
speech processing,'' \emph{IEEE JSTSP}, 2022.

\bibitem{jiwer}
Jitsi, ``JiWER: Word error rate and character error rate for automatic speech
recognition,'' \url{https://jitsi.github.io/jiwer/}.
\end{thebibliography}

\end{document}
